\documentclass[12pt,a4paper]{article}
\usepackage[utf8]{inputenc}
\usepackage[T1]{fontenc}
\usepackage[margin=2.5cm]{geometry}
\usepackage{graphicx}
\usepackage{listings}
\usepackage{xcolor}
\usepackage{caption}

\newcommand{\studentgroup}{FAF-242}
\definecolor{termbg}{rgb}{0.96,0.96,0.96}
\lstset{backgroundcolor=\color{termbg},basicstyle=\ttfamily\footnotesize,breaklines=true,frame=single}

\begin{document}
\begin{titlepage}
    \centering
    \IfFileExists{utm_logo.png}{\includegraphics[width=0.25\textwidth]{utm_logo.png}\\[0.6cm]}{\vspace*{2cm}}
    {\large \textbf{UNIVERSITATEA TEHNICĂ A MOLDOVEI}}\\[0.2cm]
    {\large Facultatea Calculatoare, Informatică și Microelectronică}\\[0.2cm]
    {\large Departamentul Ingineria Software și Automatică}\\[2.5cm]
    {\Large \textbf{LUCRAREA DE LABORATOR NR. 3}}\\[0.4cm]
    {\LARGE \textbf{Bash, scripting și filtre de text}}\\[0.4cm]
    {\large Disciplina: Sisteme de Operare}\\[2.5cm]
    \begin{minipage}{0.45\textwidth}\flushleft\textbf{Done by:}\\st. gr. \studentgroup\\Daniel-Ioan \textsc{Strunga}\end{minipage}\hfill
    \begin{minipage}{0.50\textwidth}\flushright\textbf{Checked by:}\\asist. univ. Patricia \textsc{Reițman}\\asist. univ. Marius \textsc{Dumitrașcu}\end{minipage}
    \vfill{\large Chișinău --- 2026}
\end{titlepage}
\newpage

\section{Scopul lucrării}
Familiarizarea cu linia de comandă Linux și automatizarea sarcinilor recurente prin scrierea de scripturi Bash. S-au utilizat utilitare standard (grep, cut, sort, uniq, wc) pentru a procesa, filtra și extrage statistici dintr-un fișier jurnal (\texttt{access.log}).

\section{Desfășurarea lucrării}
\subsection{Codul scripturilor (\texttt{oneliners.sh} și \texttt{report.sh})}
Am creat \texttt{oneliners.sh} pentru a răspunde cu o singură linie de comandă la fiecare cerință, utilizând pipe-uri pentru a lega output-ul unei comenzi la input-ul alteia. În \texttt{report.sh} am implementat logica de bash scripting iterativ.

\begin{figure}[h!]
\begin{lstlisting}[language=bash]
user@user-VPCEH1L0E:~/os-lab3$ ./report.sh access.log
500 requests, 100 failed
user0: 15 failures
user1: 14 failures
...
\end{lstlisting}
\caption{Figure 3.1 -- Rularea și output-ul scriptului report.sh}
\end{figure}

\subsection{Răspunsuri (Observe)}
\begin{enumerate}
    \item Cererile eșuate au fost identificate filtrând liniile care conțin cuvântul "FAIL" cu \texttt{grep -c}.
    \item Pentru paginile distincte am folosit \texttt{cut} pentru a extrage coloana paginii, urmat de \texttt{sort -u}.
    \item Utilizatorul cu cele mai multe erori a fost extras prin numărarea aparițiilor unice și sortarea lor numerică descrescătoare (\texttt{sort -rn | head -1}).
\end{enumerate}

\section{Concluzii}
Utilizarea pipe-urilor în Linux reprezintă o metodă extrem de eficientă de procesare a datelor brute. Am demonstrat că operațiuni complexe de analiză a logurilor pot fi reduse la o singură linie de cod prin combinarea filtrelor standard, eliminând nevoia unor limbaje de programare de nivel înalt pentru sarcini administrative.
\end{document}
